\ticket{15}{Модели прогнозирования на основе деревьев решений. Алгоритмы CART и C4.5: критерии поиска разбиений, параметры ограничения роста и отсечения}

\begin{examplan}
    \item Описать деревья решений как модели классификации и регрессии.
    \item Раскрыть критерии разбиения в CART и C4.5.
    \item Объяснить ограничения роста дерева.
    \item Описать отсечение и борьбу с переобучением.
\end{examplan}

\subsection*{Короткая версия для письменного ответа}

\textbf{Дерево решений} --- модель вида "если--то", которая рекурсивно разбивает пространство признаков. Внутренние узлы задают проверки признаков, ветви --- результаты проверок, листья --- прогноз: класс, вероятности классов или числовое значение.

Деревья интерпретируемы и работают с разными типами признаков, но склонны к переобучению и нестабильны к небольшим изменениям данных.

\subsection*{Построение дерева}

Дерево строится сверху вниз:
\begin{enumerate}
    \item начать с корня, содержащего всю обучающую выборку;
    \item выбрать лучший признак и порог/значение для разбиения;
    \item разделить выборку на дочерние узлы;
    \item рекурсивно повторять процесс;
    \item остановиться по критерию остановки или затем выполнить pruning.
\end{enumerate}

\subsection*{CART}

\textbf{CART} строит бинарные деревья и применяется как для классификации, так и для регрессии.

Для классификации часто используется \textbf{индекс Джини}:
\[
Gini(m)=1-\sum_{k=1}^{K}p_{mk}^2,
\]
где $p_{mk}$ --- доля объектов класса $k$ в узле $m$.

Лучшее разбиение максимизирует уменьшение нечистоты:
\[
\Delta Gini=Gini(parent)-\sum_j\frac{N_j}{N}Gini(child_j).
\]

Для регрессии CART минимизирует дисперсию или среднеквадратичную ошибку в листьях. Для числовых признаков перебираются возможные пороги, для категориальных --- разбиения категорий на группы.

\subsection*{C4.5}

\textbf{C4.5} строит деревья классификации, допускает многозначные разбиения и умеет работать с пропущенными значениями.

Основной критерий связан с энтропией:
\[
Entropy(S)=-\sum_{k=1}^{K}p_k\log_2 p_k.
\]

\textbf{Прирост информации}:
\[
Gain(S,A)=Entropy(S)-\sum_{v}\frac{|S_v|}{|S|}Entropy(S_v).
\]

Обычный прирост информации склонен выбирать признаки с большим числом значений, поэтому C4.5 использует \textbf{gain ratio} --- нормировку прироста информации на информацию разбиения.

\subsection*{Ограничение роста}

Pre-pruning останавливает рост дерева заранее. Типичные параметры:
\begin{itemize}
    \item максимальная глубина;
    \item минимальное число объектов в узле для разбиения;
    \item минимальное число объектов в листе;
    \item минимальное улучшение критерия качества;
    \item максимальное число листьев.
\end{itemize}

Ранняя остановка может быть слишком жадной и остановить дерево до полезных глубоких разбиений.

\subsection*{Отсечение}

Post-pruning сначала строит большое дерево, затем удаляет ветви, которые ухудшают обобщение.

Основные подходы:
\begin{itemize}
    \item reduced error pruning: проверка замены поддерева листом на валидационной выборке;
    \item pessimistic error pruning в C4.5: оценка ошибки с поправкой на сложность без отдельной валидации;
    \item cost-complexity pruning в CART: минимизация $R_\alpha(T)=R(T)+\alpha |L(T)|$, где $|L(T)|$ --- число листьев.
\end{itemize}

Окончательное дерево часто выбирают по кросс-валидации.

\begin{important}
    \item Дерево решений рекурсивно делит пространство признаков.
    \item CART строит бинарные деревья; для классификации часто использует Gini, для регрессии --- MSE.
    \item C4.5 использует entropy, information gain и gain ratio.
    \item Главная проблема деревьев --- переобучение.
    \item Ограничение роста и pruning уменьшают сложность дерева и улучшают обобщение.
\end{important}
